\documentclass[11pt,a4paper]{article}
\usepackage[utf8]{inputenc}
\usepackage[T1]{fontenc}
\usepackage{amsmath,amssymb,amsthm,mathtools}
\usepackage{bm}
\usepackage{hyperref}
\usepackage[capitalise]{cleveref}
\usepackage{microtype}
\hypersetup{colorlinks=true, linkcolor=blue, citecolor=blue}

\DeclarePairedDelimiter{\abs}{\lvert}{\rvert}
\DeclarePairedDelimiter\norm{\lVert}{\rVert}
\DeclarePairedDelimiter{\ceil}{\lceil}{\rceil}
\DeclareMathOperator*{\argmin}{arg\,min}
\DeclareMathOperator{\Tr}{Tr}
\newcommand{\R}{\mathbb{R}}
\newcommand{\inner}[2]{\left\langle #1, #2 \right\rangle}

\theoremstyle{plain}
\newtheorem{theorem}{Theorem}[section]
\newtheorem{lemma}[theorem]{Lemma}
\theoremstyle{definition}
\newtheorem{definition}[theorem]{Definition}
\theoremstyle{remark}
\newtheorem*{remark}{Remark}

\title{Convex Optimization Notes\thanks{Draft version.}}
\author{A. Student}
\date{\today}

\begin{document}
\maketitle

\begin{abstract}
We collect standard facts about convex functions and gradient descent.
\end{abstract}

\section{Preliminaries}\label{sec:prelim}

\begin{definition}[Convexity]\label{def:convex}
A function $f\colon \R^n \to \R$ is \emph{convex} if for all $\bm{x}, \bm{y}$ and $\lambda \in [0,1]$,
\[
  f(\lambda \bm{x} + (1-\lambda)\bm{y}) \le \lambda f(\bm{x}) + (1-\lambda) f(\bm{y}).
\]
\end{definition}

Let $L \coloneqq \sup_{x} \norm{\nabla^2 f(x)}$ and define $x^\star \coloneqq \argmin_{x \in \R^n} f(x)$.
We write $\abs*{x}$ for the absolute value and $\ceil{x}$ for the ceiling.

\begin{lemma}\label{lem:descent}
If $f$ is $L$-smooth then
\begin{equation}\label{eq:descent}
  f(y) \le f(x) + \inner{\nabla f(x)}{y - x} + \frac{L}{2}\norm{y - x}^2 .
\end{equation}
\end{lemma}

\begin{proof}[Proof of \cref{lem:descent}]
By Taylor's theorem,
\begin{align*}
  f(y) &= f(x) + \int_0^1 \inner{\nabla f(x + t(y-x))}{y-x}\,dt \\
       &\le f(x) + \inner{\nabla f(x)}{y-x} + \frac{L}{2}\norm{y-x}^2. \qedhere
\end{align*}
\end{proof}

\begin{theorem}\label{thm:rate}
Gradient descent with step $1/L$ satisfies
\begin{equation*}
  f(x_k) - f(x^\star) \le \frac{L \norm{x_0 - x^\star}^2}{2k}.
\end{equation*}
\end{theorem}

See \cref{thm:rate,lem:descent}, \Cref{eq:descent}, and \crefrange{lem:descent}{thm:rate}.
Also \autoref{sec:prelim} and \nameref{sec:prelim}.

\begin{remark}
The rate in \cref{thm:rate} is not optimal.
\end{remark}

\section{Piecewise definitions}

\begin{equation}
  \operatorname{sgn}(x) = \begin{dcases}
    -1 & x < 0, \\
    \phantom{-}0 & x = 0, \\
    \frac{x}{\abs{x}} & \text{otherwise}.
  \end{dcases}
\end{equation}

\begin{equation*}
  g(x) = \begin{cases*}
    x^2 & if $x \ge 0$ \\
    0 & otherwise
  \end{cases*}
\end{equation*}

\begin{align}
  a &= b + c \\
  \shortintertext{and therefore}
  d &= \mathclap{\sum_{i=1}^{n}} e_i
\end{align}

We also use $\prescript{14}{6}{\mathrm{C}}$ and $\Tr(A^\top A) = \norm{A}_F^2$.
\begin{equation}
  \begin{pmatrix*}[r] -1 & 2 \\ 3 & -4 \end{pmatrix*}
  \begin{bmatrix} 1 \\ 0 \end{bmatrix}
  = \begin{pmatrix} -1 \\ 3 \end{pmatrix}
\end{equation}

\end{document}
